\scanpage{59}
% archive sheet number: 58
\unsure{section} $s$ de $\underline{S}^{n}(\underline{E})$ s'identifie à une \unsure{section} de $\underline{L}_{\underline{E}}$ sur $X=P(\underline{E})$, et \deleted{\unsure{a}} \inserted{définit} donc un \uline{diviseur} $\operatorname{div}(s)\deleted{\unsure{(s)}}$ sur $X$; la \unsure{condition} que $s$ [sur $Y$] ne \uline{s'annule pas}, signifie que ce diviseur induit bien un diviseur sur chaque fibre, \deleted{\unsure{D}} i.e. que $s$ ne s'annule pas.

D'ailleurs \deleted{\unsure{cette correspondance}} $\operatorname{div}(s)=\operatorname{div}(t)$ \deleted{\unsure{si et}} si et seulement si on \unsure{peut} trouver $\rho\in\Gamma(X,\underline{O}_X^*)$ t.q. $t=\rho s$, or $\rho$ s'identifie à un élément de $\Gamma(Y,\underline{O}_Y^*)$, (car $f_*(\unsure{\mathscr{O}|\underline{O}}_X)=\underline{O}_Y$). Donc $\operatorname{div}(s)=\operatorname{div}(t)\Leftrightarrow s$ et $t$ définissent la $=$ section de $\underline{S}^{n}(\underline{E})^*/\underline{O}_Y^*$.

De plus, $s\unsure{-|\to}\operatorname{div}(s)$ est compatible aux localisations sur $S$. D'où : \uline{on a une application injective} de \deleted{$\unsure{\Phi|\varphi}$} $\underline{\Phi}^{(n)}(X)(Y)=\Gamma(Y,\underline{S}^{n}(\underline{E})^*/\underline{O}_Y^*)$ \uline{dans l'ens des diviseurs loc. principaux \inserted{positifs} sur $X$ qui induisent \deleted{un diviseur} sur chaque fibre un diviseur, \inserted{qui est} de degré $n$}. Appelons de tels diviseurs sur $X$ des ``diviseurs transversaux de degré $n$''.

Je dis que l'application précédente est \uline{surjective}, i.e. que tout diviseur transversal \inserted{$D$} de degré $n$ provient d'une section de $\Phi^{(n)}(X)$. En effet, soit $\underline{L}=\underline{L}(D)$ le faisceau inversible défini par
